\documentclass[11pt]{article}

\usepackage[margin=1in]{geometry}
\usepackage{amsmath,amssymb}
\usepackage{enumitem}
\usepackage[dvipsnames]{xcolor}
\usepackage[colorlinks=true,linkcolor=blue,citecolor=blue,urlcolor=blue]{hyperref}
\usepackage{setspace}
\usepackage[T1]{fontenc}
\usepackage{lmodern}

\setlength{\parindent}{0pt}
\setlength{\parskip}{0.6em}

% =========================================================
% Commands
% =========================================================

\newcommand{\reply}{
\vspace{0.25em}
{\color{ForestGreen}\textbf{Reply:}}
}

\newcommand{\changes}{
\vspace{0.25em}
{\color{BrickRed}\textbf{Changes:}}
}

\newenvironment{reviewercomment}
{
\vspace{0.5em}
\begin{quote}
\color{NavyBlue}
}
{
\end{quote}
}

\newcommand{\TFT}{\emph{TFT}}
\newcommand{\AXL}{Axelrod-Python}

% =========================================================
% Document
% =========================================================

\begin{document}

We thank the editor for handling our submission, and we are grateful to the
reviewers for their careful reading of the manuscript and their constructive
feedback. Their comments have strengthened the paper and helped us improve both
its clarity and overall presentation. In this document, we provide detailed
responses to each comment.

% =========================================================
% Reviewer 1
% =========================================================

\section*{Reviewer 1}

\begin{reviewercomment}
The reproducibility of results has become an important issue across fields. It
remains an open question whether Axelrod's second tournament, which contributed
significantly to research on direct reciprocity, can be reproduced. In this
manuscript, the authors adapt the surviving Fortran implementations to modern
compilers and integrate them into the Axelrod-Python framework. Despite the
stochasticity of the game setting, the reproduced results are consistent with
the original ones. The authors also provide a detailed analysis of the
tournament results. In addition, they explore additional tournaments by
introducing extra strategies and incorporating noise to test the performance.
\end{reviewercomment}

\reply

We sincerely thank the reviewer for their positive assessment of our work and
for recognizing its contributions.

\begin{reviewercomment}
I have the following comments and questions about the manuscript:

1. Do the original 63 submitted strategies include classical benchmark
strategies such as AllC, AllD, or WSLS? If so, it would be better to report
their rankings.
\end{reviewercomment}

\reply

The original set of 63 submitted
strategies did not include classical benchmark strategies such as AllC, AllD,
or WSLS. The only strategies that would now be considered classical were Tit
For Tat and Random. Motivated by this comment, and the related point raised
below, we carried out an additional analysis to examine what happens when such
benchmark strategies are included.

\changes

We have added a sentence clarifying that classical benchmark strategies were not
included in the original tournament (in the ``Reproducing Axelrod's Second
Tournament'' section, under ``Running the tournament'').

\begin{reviewercomment}
2. Regarding the Extra Invitations section, the win rate of TFT exhibits
non-monotonic fluctuations as the number of additional strategies increases. Is
this behavior primarily driven by the distributional properties of the 209
Axelrod-Python strategies? 

Moreover, this section seems to be motivated by
computational exploration. Therefore, in addition to showing that TFT and k42r
remain relatively robust when additional strategies are introduced, I suggest
presenting results that incorporate classical strategies and examining how they
affect the performance of TFT and k42r.
\end{reviewercomment}

\reply

The reviewer raises two related points. First, they ask about the
non-monotonic behavior of the winning probability of \TFT{} as the number of
additional strategies increases. Second, they suggest incorporating classical
benchmark strategies into the analysis in order to better understand how the
performance of \TFT{} and \texttt{k42r} changes in extended tournaments.

To this end, we first consider the effect of introducing several classical
strategies into the original tournament. We find that, among the strategies we
consider, only Anti-\TFT{} is capable of dethroning \TFT{}. This occurs because
Anti-\TFT{} is able to exploit \TFT{}, lowering its payoff relative to other
strategies. At the same time, \texttt{k42r} is able to exploit Anti-\TFT{},
achieving payoffs higher than those obtained through mutual cooperation and
ultimately winning the tournament.

We then extend this analysis by considering all possible combinations of these
classical strategies. In this setting, we again observe a non-monotonic
relationship between the number of added strategies and the winning probability
of \TFT{}. However, unlike in the large-scale tournament analysis, here we can
explicitly identify the mechanism driving this behavior.

In particular, whenever Anti-\TFT{} is included in the tournament, \TFT{}
ranks second while \texttt{k42r} ranks first. As a result, the winning
probability of \TFT{} is determined entirely by the number of subsets that
contain Anti-\TFT{}. For example, when two additional strategies are
introduced, there are \(36\) possible combinations, of which \(7\) include
Anti-\TFT{}. These are precisely the tournaments in which \TFT{} fails to win.
The same combinatorial symmetry appears for subset sizes seven and two, six and
three, and so forth, producing the approximately symmetric non-monotonic
pattern observed in the manuscript.

We believe this analysis sheds additional light on the large-scale tournament
results. In particular, it suggests that the observed non-monotonic behavior
is not simply a property of the distribution of strategies, but can emerge from
the combinatorial interaction structure induced by specific subsets of
opponents. At the scale of the full \AXL{} tournament, however, the number of
possible interaction structures becomes extremely large, making it difficult to
fully characterize the precise mechanisms responsible for the observed
fluctuations.

\changes

We now begin the revisiting section by considering the case where classical
benchmark strategies are included in the tournament. We first examine the
effect of introducing these strategies individually and simultaneously
(Table~1), and then extend the analysis to all possible combinations of the
classical strategies for subset sizes ranging from 2 to 8 (Table~2).

We also include a new figure (Fig.~4) showing that, although the winning
probability of \TFT{} decreases for intermediate subset sizes, \TFT{} remains
consistently associated with higher average payoffs across tournaments.
Interestingly, the reduction in \TFT{}'s effectiveness coincides with the
largest variation in tournament cooperation levels, suggesting that \TFT{}
struggles most in environments with highly variable interaction structures.

Finally, we now discuss the non-monotonic behavior of \TFT{} in the
\AXL{} tournament setting. While the classical strategy analysis allows us to
identify a concrete combinatorial mechanism driving this behavior, a complete
understanding of the large-scale tournament results remains difficult due to
the combinatorial explosion of possible interaction structures and opponent
subsets.

\begin{reviewercomment}
3. In a tournament, a strategy's performance is evaluated against a specific
distribution of opponents. In Axelrod's second tournament, this distribution is
defined by the other 62 submitted strategies; in the large-scale experiments,
the Axelrod-Python strategy set is also included. As both distributions are
subjective choices, is it possible to construct an "objective" distribution of
opponents and obtain a more "objective" performance ranking accordingly?
\end{reviewercomment}

\reply

In practice, the set of participating strategies determines the
types of interactions that can occur within the tournament. As demonstrated by
our extra invitation experiments, modifying the tournament composition can
substantially change the resulting rankings.

However, this also raises the question of what exactly we mean by a
``distribution of opponents'' and what it would mean for such a distribution
to be ``objective''. One possible interpretation is that the opponent set
induces a distribution over interaction types. For example, in the
Supplementary Information, we plot the distributions of cooperation rates and
payoffs across pairwise interactions for both the original tournament and the
\AXL{} tournament. These distributions are highly skewed, indicating
that the tournaments do not sample interactions uniformly across behavioural
space.

This naturally raises further questions. Should an ``objective'' tournament
sample interactions uniformly across behavioural space? Should it instead
balance different levels of cooperation, reciprocity, exploitability, or payoff
asymmetry? More generally, what properties should define a representative
interaction environment for evaluating repeated game strategies?

We believe this is an important open question.

\changes

We have added a discussion of this point in the manuscript. We also include new
figures in the Supplementary Information (S7) illustrating the distributions of
interaction cooperation rates and payoffs in both the original and expanded
tournaments, highlighting their highly skewed structure.

\begin{reviewercomment}
Overall, this manuscript successfully revisits and extends a foundational work
in evolutionary game theory and serves as a valuable case study in computational
reproducibility. I recommend the manuscript for publication after minor
revisions.
\end{reviewercomment}

\reply

We thank the reviewer.

% =========================================================
% Reviewer 2
% =========================================================

\section*{Reviewer 2}

\begin{reviewercomment}
The authors reproduced one of the famous tournaments conducted by Robert Axelrod
about half a century ago. Axelrod's main findings were reconfirmed. The authors
also tested some variations, e.g., by using an extended set of strategies and
adding noise. This is an interesting paper. Reproducibility is an important part
of academic integrity, and the preservation of this historical event is a
meaningful activity. As long as this lies within the scope of this journal, I
would like to recommend accepting this manuscript for publication in
communications ai \& computing.
\end{reviewercomment}

\reply

We thank the reviewer for their careful and thoughtful reading of our
manuscript, as well as for their encouraging remarks about its contributions.

This work falls naturally within social computing and computational social
science, as it studies how interaction environments shape cooperation and
collective outcomes in repeated social dilemmas. By revisiting and extending
Axelrod's tournaments through large-scale computational experiments, we examine
how strategic behaviour, reciprocity, and opponent composition influence
tournament outcomes and the emergence of cooperation.

We believe Communications AI \& Computing is an excellent fit for this work
because the manuscript combines large-scale computational modelling,
reproducible software infrastructure, historical computational reproducibility,
and the analysis of social interaction systems within the context of social
computing.

% =========================================================
% Reviewer 3
% =========================================================

\section*{Reviewer 3}

\begin{reviewercomment}
This outstanding paper faithfully revives Axelrod's second IPD tournament by
compiling the original 63 Fortran submissions (with only compiler fixes) and
integrating them into Axelrod-Python. It reproduces the core result, TFT (k92r)
wins, while uncovering a bug in Champion (k61r) and showing k42r wins ~31.5\% of
replications. Extensions to noisy tournaments (1\%) and massive 272-strategy
tournaments (with Axelrod-Python + Stewart \& Plotkin) reveal TFT's robustness in
cooperative settings but its vulnerability in diverse/noisy environments. 
\end{reviewercomment}

\reply

We thank the reviewer for their detailed and positive assessment of our work.
We are grateful for their recognition of the contributions of our paper, as
well as for their constructive feedback on how to further strengthen the
manuscript.

\begin{reviewercomment}
None that threaten validity, but two points merit clarification or strengthening
for publication: (1) The authors infer the fifth length as 156 from the reported
average of 151 moves. This is reasonable but not 100\% certain (Axelrod's
original announcement may have used slightly different probabilistic sampling).
Perhaps a short paragraph acknowledging this (and perhaps providing sensitivity
results with alternative length sets, e.g., {63, 77, 151, 200, 308} or full
probabilistic sampling) would be useful. 
\end{reviewercomment}

\reply

Our choice of a match length of 156 follows directly from Axelrod's reported
average match length of 151 moves under the five tournament lengths used in the
second tournament. However, the incomplete description of the stopping rule in
the original paper introduces some uncertainty regarding the precise
implementation of the tournament.

To address this point, we now additionally perform a sensitivity analysis in
which matches end probabilistically rather than after a fixed number of turns.
Specifically, we consider a continuation probability corresponding to an
expected match length of 151 turns.

Overall, we find that the main conclusions remain unchanged under probabilistic
ending conditions. In particular, the ranking of strategies remains consistent
with the original analysis, and \TFT{} continues to perform robustly.

\changes
We have included a new section in the Supplementary Information (S8)
describing the sensitivity analysis with probabilistic match ending. In this
analysis, we reproduce the original tournament using a true probabilistic
stopping rule rather than fixed sampled match lengths. We use a continuation
probability corresponding to an expected match length of 151 turns and repeat
the full tournament 1000 times for each of five random seeds. We now
discuss this analysis in the Conclusion section.

\begin{reviewercomment}
(2) The 78+ million tournaments are subsampled from the full interaction matrix.
Please clarify the subsampling method (random draws of the added strategies?)
and report confidence intervals or standard errors on the proportions
(especially the non-monotonic TFT behaviour).
\end{reviewercomment}

\reply

We have substantially
revised the extra invitations section in response to Reviewer~1's comments,
which we believe now makes the construction of these tournaments much clearer.

For each subset size, we consider all possible combinations of strategies drawn
from the given strategy set. Thus, the tournaments are not generated through
random subsampling of additional strategies; rather, we exhaustively enumerate
all possible subsets up to the specified subset size. As a result, the reported
winning probabilities correspond to deterministic frequencies over the full
combinatorial space of tournaments considered, rather than estimates obtained
from a random sample.

We agree, however, that the original text did not explain this procedure
clearly enough. We have therefore clarified the construction of the tournaments
and expanded the discussion of the observed non-monotonic behavior of
\TFT{} in the revised manuscript.

\changes

We have clarified the construction of the extra invitation tournaments in the
manuscript.



\begin{reviewercomment}
Minor issues:

(1) ``k42r ranked only 3th'' $\rightarrow$ ``3rd''.

(2) ``78,760,605 tournaments'', please double-check arithmetic.
\end{reviewercomment}

\reply

We have corrected the
typo (``3th'' to ``3rd''). We also clarified that the reported
number of tournaments refers to the sum of all combinations obtained by adding
between one and four strategies. The correct total is
78{,}760{,}569 tournaments.

\changes

Corrected in the revised manuscript.

\begin{reviewercomment}
This is a high-quality, timely, and impactful contribution. I strongly
recommend acceptance after these light revisions. The paper will be widely
read and cited, both for its scientific findings and as a case study in
digital preservation of historical computational science.
\end{reviewercomment}

\reply

We thank the reviewer.

\end{document}